% (c) Nikita Lisitsa, lisyarus@gmail.com, 2024

\documentclass[10pt]{beamer}

\usepackage[T2A]{fontenc}
\usepackage[russian]{babel}
\usepackage{minted}

\usepackage{graphicx}
\graphicspath{ {./images/} }

\usepackage{adjustbox}

\usepackage{color}
\usepackage{soul}

\usepackage{hyperref}

\usetheme{metropolis}

\definecolor{red}{rgb}{1,0,0}
\definecolor{green}{rgb}{0,0.5,0}
\definecolor{blue}{rgb}{0,0,1}
\definecolor{codebg}{RGB}{29,35,49}
\definecolor{lightbg}{RGB}{253,246,227}
\setminted{fontsize=\footnotesize}

\makeatletter
\newcommand{\slideimage}[1]{
  \begin{figure}
    \begin{adjustbox}{width=\textwidth, totalheight=\textheight-2\baselineskip-2\baselineskip,keepaspectratio}
      \includegraphics{#1}
    \end{adjustbox}
  \end{figure}
}
\makeatother

\title{Фотореалистичный рендеринг\quad\quad\quad\quad\quad\quad \textit{(aka raytracing)}}
\subtitle{Лекция 8: Текстуры, environment maps, GPU raytracing}
\date{2024}

\setbeamertemplate{footline}[frame number]

\begin{document}

\frame{\titlepage}

\begin{frame}
\frametitle{Зачем нужны текстуры?}
\begin{itemize}
\item Дешевле детализация -- добавление геометрии не бесплатное
\pause
\item Возможность интерполировать параметры материала между вершинами
\pause
\item Отделение геометрии от детализации удобнее художникам
\pause
\item Можно поддержать environment maps
\end{itemize}
\end{frame}

\begin{frame}
\frametitle{Какие бывают текстуры?}
\begin{itemize}
\item Любой параметр, влияющий на BRDF в точке, можно записать в текстуру:
\pause
\begin{itemize}
\item Цвет
\pause
\item Вектор нормали (\textit{normal mapping})
\pause
\item Параметры материала, e.g. metallic и roughness (\textit{material mapping})
\pause
\item Цвет излучения
\end{itemize}
\pause
\item Environment map -- текстура `окружения' сцены
\end{itemize}
\end{frame}

\begin{frame}
\frametitle{Normal mapping}
\slideimage{sphere_normal.png}
\end{frame}

\begin{frame}
\frametitle{Normal mapping}
\slideimage{sphere-specular.png}
\end{frame}

\begin{frame}
\frametitle{Normal mapping}
\slideimage{sphere-metallic.png}
\end{frame}

\begin{frame}
\frametitle{Material mapping}
\slideimage{sphere_material.png}
\end{frame}

\begin{frame}
\frametitle{Material mapping}
\slideimage{sphere-roughness.png}
\end{frame}

\begin{frame}
\frametitle{Emission texture}
\slideimage{sphere_emission.png}
\end{frame}

\begin{frame}
\frametitle{Emission texture}
\slideimage{sphere-emissive.png}
\end{frame}

\begin{frame}
\frametitle{Как текстуры описываются в glTF?}
\usemintedstyle{solarized-light}
\begin{itemize}
\item Текстуры описаны в массиве \mintinline{cpp}|textures|
\pause
\item Материал может ссылаться (по индексу) на текстуры для различных параметров:
\pause
\begin{itemize}
\item \mintinline{cpp}|material.normalTexture|
\item \mintinline{cpp}|material.occlusionTexture|
\item \mintinline{cpp}|material.emissiveTexture|
\item \mintinline{cpp}|material.pbrMetallicRoughness.baseColorTexture|
\item \mintinline{cpp}|material.pbrMetallicRoughness.metallicRoughnessTexture|
\end{itemize}
\end{itemize}
\end{frame}

\begin{frame}
\frametitle{Как текстуры описываются в glTF?}
\usemintedstyle{solarized-light}
\begin{itemize}
\item Сама текстура ссылается на два других объекта:
\pause
\begin{itemize}
\item \mintinline{cpp}|texture.source| -- индекс объекта в массиве \mintinline{cpp}|images|, описывающего расположение пикселей текстуры
\pause
\item \mintinline{cpp}|texture.sampler| -- индекс объекта в массиве \mintinline{cpp}|samplers|, описывающего параметры чтения пикселей текстуры
\end{itemize}
\pause
\item \mintinline{cpp}|texture.source| может ссылаться на файл, на URL, на buffer view, etc -- мы будем поддерживать только файлы
\end{itemize}
\end{frame}

\begin{frame}
\frametitle{glTF sampler}
\usemintedstyle{solarized-light}
\begin{itemize}
\item Sampler описывает типичные настройки текстуры при realtime рендеринге:
\pause
\begin{itemize}
\item Minification и magnification фильтры
\pause
\item Wrapping modes (repeat, mirrored repeat, etc)
\end{itemize}
\pause
\item В path tracing'е разделение на min и mag фильтры не имеет большого смысла
\pause
\item Мы будем использовать \underline{всегда} LINEAR фильтр и REPEAT wrapping mode
\end{itemize}
\end{frame}

\begin{frame}
\frametitle{Wrapping modes}
\slideimage{wrapping.png}
\end{frame}

\begin{frame}
\frametitle{Текстурные координаты}
\usemintedstyle{solarized-light}
\begin{itemize}
\item Для обращения к текстуре нам нужно знать соответствие вершин модели точкам на изображении, т.е. знать \textit{текстурные координаты}
\pause
\item Они заданы так же, как другие атрибуты вершин
\pause
\item glTF поддерживает несколько разных атрибутов текстурных координат, использующихся для разных текстур
\pause
\item Мы будем поддерживать только один такой атрибут: \mintinline{cpp}|TEXCOORD_0|
\pause
\item Будем считать, что это всегда \mintinline{cpp}|VEC2| floating-point
\end{itemize}
\end{frame}

\begin{frame}
\frametitle{Формат текстур}
\usemintedstyle{solarized-light}
\begin{itemize}
\item glTF чётко описывает, как должны интерпретироваться пиксели различных текстур:
\pause
\begin{itemize}
\item \mintinline{cpp}|baseColorTexture| хранится в sRGB-формате, т.е. при чтении её пиксели нужно возвести в степень \begin{math}2.2\end{math}
\pause
\item \mintinline{cpp}|emissiveTexture| хранится в sRGB-формате
\pause
\item \mintinline{cpp}|metallicRoughnessTexture| хранится в линейном формате, где зелёный канал это roughness, а синий -- metallic
\pause
\item \mintinline{cpp}|normalTexture| хранится в линейном формате, с преобразованием \begin{math}[-1,1]\rightarrow[0,1]\end{math}
\end{itemize}
\end{itemize}
\end{frame}

\begin{frame}
\frametitle{Формат текстур}
\usemintedstyle{solarized-light}
\begin{itemize}
\item Если какая-то текстура не задана, она считается полностью белой
|pause
\item \mintinline{cpp}|baseColorFactor|, \mintinline{cpp}|emissiveFactor|, \mintinline{cpp}|roughnessFactor| и \mintinline{cpp}|metallicFactor| понимаются как \textit{коэффициенты}, на которые домножается значение из текстуры (после преобразования из sRGB и интерполяции)
\end{itemize}
\end{frame}

\begin{frame}
\frametitle{Семплинг текстуры}
\usemintedstyle{solarized-light}
\begin{itemize}
\item Удобно из функции пересечения луча с треугольником вернуть также проинтерполированные текстурные координаты
\pause
\item Удобно завести отдельную функцию, возвращающую цвет из текстуры по текстурным координатам
\pause
\item Сначала нужно произвести \textit{wrapping}: фактически, взять каждую координату по модулю 1 (учтите, что они могут быть отрицательными!)
\pause
\item После текстурные координаты нужно перевести из диапазона \begin{math}[0,1]\times [0,1]\end{math} в диапазон \begin{math}[0,\text{width}]\times[0,\text{height}]\end{math}
\pause
\item По координате \mintinline{cpp}|tx| можно вычислить координаты двух соседних пикселей \mintinline{cpp}|ix = floor(tx)| и \mintinline{cpp}|ix + 1|, и коэффициент интерполяции между ними \mintinline{cpp}|dx = tx - floor(tx)|
\pause
\item Аналогично по Y
\end{itemize}
\end{frame}

\begin{frame}
\frametitle{Семплинг текстуры}
\usemintedstyle{solarized-light}
\begin{itemize}
\item Мы получили координаты четырёх соседних пикселей (\mintinline{cpp}|ix + 1| нужно брать по модулю \mintinline{cpp}|width|, и аналогично для Y)
\pause
\item Читаем четыре значения пикселей, переводим из байтов \begin{math}[0,255]\end{math} в нормированные единицы \begin{math}[0,1]\end{math}
\pause
\item Если наша текстура -- sRGB, возводим их покомпонентно в степень \begin{math}2.2\end{math} (альфа-компоненту возводить не нужно)
\pause
\item Делаем билинейную интерполяцию полученных значений с коэффициентами \mintinline{cpp}|dx,dy|
\pause
\item Полученное значение и будет семплом нашей текстуры
\end{itemize}
\end{frame}

\begin{frame}
\frametitle{Base color texture}
\usemintedstyle{solarized-light}
\begin{itemize}
\item Для поддержки baseColorTexture (\textit{альбедо}) достаточно заменить цвет в вычислении BRDF с \mintinline{cpp}|baseColorFactor| на \mintinline{cpp}|baseColorFactor * sample(baseColorTexture, texcoord).rgb|
\end{itemize}
\end{frame}

\begin{frame}
\frametitle{Emissive texture}
\usemintedstyle{solarized-light}
\begin{itemize}
\item Для поддержки emissiveTexture достаточно заменить цвет излучения с \mintinline{cpp}|emissiveFactor| на \mintinline{cpp}|emissiveFactor * sample(emissiveTexture, texcoord).rgb|
\pause
\item Не забудьте, что мы поддерживаем расширение \mintinline{cpp}|KHR_materials_emissive_strength|, которое даёт дополнительныый множитель для \mintinline{cpp}|emissiveFactor|
\end{itemize}
\end{frame}

\begin{frame}
\frametitle{Metallic roughness texture}
\usemintedstyle{solarized-light}
\begin{itemize}
\item Для поддержки metallicRoughnessTexture нужно:
\pause
\begin{itemize}
\item Получить семпл текстуры \mintinline{cpp}|sample = sample(metallicRoughnessTexture, texcoord).rgb|
\pause
\item Вычислить roughness как \mintinline{cpp}|roughnessFactor * sample.g|
\pause
\item Вычислить metallic как \mintinline{cpp}|metallicFactor * sample.b|
\pause
\item Использовать их в рассчётах BRDF
\end{itemize}
\end{itemize}
\end{frame}

\begin{frame}
\frametitle{Normal texture}
\usemintedstyle{solarized-light}
\begin{itemize}
\item Normal texture -- самая неприятная
\pause
\item Нормали из этой текстуры задаются в локальной системе координат поверхности, ось Z которой совпадает с shading normal, а ось X -- с \textit{tangent} вектором
\pause
\item Tangent тоже задан как атрибут вершин \mintinline{cpp}|TANGENT|, мы будем считать, что он всегда \mintinline{cpp}|VEC4| floating-point
\pause
\item Первые три координаты \mintinline{cpp}|TANGENT| задают координаты локальной оси X в системе координат объекта
\pause
\item Четвёртая координата \mintinline{cpp}|TANGENT| равна 1, если локальная система координат -- правая, или -1, если левая
\end{itemize}
\end{frame}

\begin{frame}[fragile]
\frametitle{Normal texture}
\usemintedstyle{solarized-light}
\begin{itemize}
\item Tangent вектор нужно преобразовывать в мировую систему координат так же, как координаты самих вершин
\pause
\item Четвёртую координату при этом нужно оставить как есть!
\pause
\item Зная normal и tangent, локальную систему координат можно восстановить как
\begin{minted}[bgcolor=lightbg]{cpp}
local_X = tangent.xyz
local_Z = normal
local_Y = cross(local_Z, local_X) * tangent.w
\end{minted}
\end{itemize}
\end{frame}

\begin{frame}[fragile]
\frametitle{Normal texture}
\usemintedstyle{solarized-light}
\begin{itemize}
\item Значение из normalTexture нужно перевести в диапазон \begin{math}[-1,1]\end{math} и перевести из локальной системы координат в глобальную
\begin{minted}[bgcolor=lightbg]{cpp}
sample = sample(normalTexture, texcoord)
local_normal = sample * 2.0 - vec3(1.0)
normal = local_X * local_normal.x
       + local_Y * local_normal.y
       + local_Z * local_normal.z
normal = normalize(normal)
\end{minted}
\pause
\item Эту нормаль нужно использовать, как \textit{shading normal}
\end{itemize}
\end{frame}

\begin{frame}
\frametitle{Sponza без текстур}
\slideimage{sponza-white.png}
\end{frame}

\begin{frame}
\frametitle{Sponza, только альбедо}
\slideimage{sponza-textured-albedo-only.png}
\end{frame}

\begin{frame}
\frametitle{Sponza, все текстуры}
\slideimage{sponza-textured-full.png}
\end{frame}

\begin{frame}
\frametitle{Комната}
\slideimage{room.png}
\end{frame}

\begin{frame}
\frametitle{Šibenik cathedral}
\slideimage{sibenik-day.png}
\end{frame}

\begin{frame}
\frametitle{Šibenik cathedral}
\slideimage{sibenik-night.png}
\end{frame}

\begin{frame}[fragile]
\frametitle{Environment maps}
\usemintedstyle{solarized-light}
\begin{itemize}
\item glTF не поддерживает environment maps, только с помощью довольно сложного расширения \mintinline{cpp}|EXT_lights_image_based|
\pause
\item Environment maps очень удобны для визуализации сцен, в которых нет собственных источников света
\pause
\item Environment maps почти всегда хранятся в HDR формате, т.е. их пиксели не ограничены диапазоном \begin{math}[0,1]\end{math}, и почти всегда используют географическую проекцию:
\begin{minted}[bgcolor=lightbg]{cpp}
texcoord.x = 0.5 + 0.5 * atan2(dir.z, dir.x) / pi
texcoord.y = 0.5 - 0.5 * asin(dir.y)
\end{minted}
\pause
\item Можно использовать environment map как значение `фона', если луч не пересёк никакую другую геометрию
\pause
\item На сайте \href{https://hdri-haven.com}{\texttt{hdri-haven.com}} можно скачать очень много HDR environment maps
\pause
\item \textbf{\alert{N.B.}}: Значения в HDR бывают \textbf{очень} большими (например, \begin{math}10^6\end{math}), что может привести к артефактам при рендеринге
\end{itemize}
\end{frame}

\begin{frame}[fragile]
\frametitle{GPU raytracing}
\usemintedstyle{solarized-light}
\begin{itemize}
\item Используемый нами алгоритм -- Unidirectional Monte-Carlo Path Tracing -- относится к категории \textit{massively parallel}
\pause
\item Логично попробовать его портировать на GPU!
\pause
\item Например, нарисуем полноэкранный прямоугольник, и фрагментный шейдер будет для каждого пикселя вычислять его цвет, усредняя запрошенное количество семплов
\end{itemize}
\end{frame}

\begin{frame}[fragile]
\frametitle{GPU raytracing}
\usemintedstyle{solarized-light}
\begin{itemize}
\item На GPU мы столкнёмся с некоторыми проблемами:
\pause
\begin{itemize}
\item Много семплов генерировать дорого -- драйвер может решить, что GPU зависла, и перезапустить её, убив нашу программу
\pause
\item Нужно как-то передать информацию о сцене на GPU
\pause
\item Языки шейдеров довольно примитивные -- придётся отказаться от иерархий классов в пользу индексов, битовых полей, и т.п.
\pause
\item Основные компоненты нашего алгоритма -- обход BVH и сам path tracing -- рекурсивные, что тоже не поддерживается шейдерами
\end{itemize}
\end{itemize}
\end{frame}

\begin{frame}[fragile]
\frametitle{Генерация семплов}
\usemintedstyle{solarized-light}
\begin{itemize}
\item Проблему с зависанием проще всего решить, рисуя по одному семплу за раз, и усредняя их встроенным blending'ом
\pause
\item На шаге N мы хотим взять уже записанное усреднённое значение с коэффициентом \begin{math}\frac{N-1}{N}\end{math} и новый семпл с коэффициентом \begin{math}\frac{1}{N}\end{math} -- это и будет значение альфа-канала для блендинга
\end{itemize}
\end{frame}

\begin{frame}[fragile]
\frametitle{Информация о сцене}
\usemintedstyle{solarized-light}
\begin{itemize}
\item Все современные графические API поддерживают передачу больших буферов на GPU
\pause
\item Например, в OpenGL это можно сделать через SSBO (shader storage buffer object), или через buffer textures
\pause
\item В Vulkan это storage buffers
\pause
\item В DirectX это structured buffers
\pause
\item И т.п.
\end{itemize}
\end{frame}

\begin{frame}[fragile]
\frametitle{Рекурсия}
\usemintedstyle{solarized-light}
\begin{itemize}
\item Рекурсию в обходе BVH легко заменить на локальный стек небольшого размера (напр. 64)
\pause
\item При спуске в дочерние вершины мы делаем одну из них текущей вершиной, а вторую кладём в стек
\pause
\item Когда текущая вершина обработана, достаём вершину из стека
\end{itemize}
\end{frame}

\begin{frame}[fragile]
\frametitle{Рекурсия}
\usemintedstyle{solarized-light}
\begin{itemize}
\item Рекурсию в самом алгоритме path tracing'а легко заменить на цикл
\pause
\item В конце одного шага path tracing'а код всегда выглядит как-то так:
\begin{minted}[bgcolor=lightbg]{cpp}
return emission + brdf * ray_color(new_ray)
\end{minted}
\pause
\item Нам достаточно запоминать суммарное количество излучения уже посещённых объектов, и множитель, на который нужно домножить следующий луч
\end{itemize}
\end{frame}

\begin{frame}[fragile]
\frametitle{Рекурсия}
\usemintedstyle{solarized-light}
\begin{minted}[bgcolor=lightbg]{cpp}
result = vec3(0.0)
factor = vec3(1.0)

for (bounce in [0..max))
    ...
    result += emission * factor
    factor *= brdf
    ray = reflected_ray

return result;
\end{minted}
\end{frame}

\begin{frame}[fragile]
\frametitle{Mega kernel}
\usemintedstyle{solarized-light}
\begin{itemize}
\item Такой подход называется \textit{mega kernel}: один большой шейдер делаёт всё
\pause
\item Так мы уже получим рабочий path tracer, но его производительность будет не очень хорошей из-за \textit{divergence}
\pause
\item \textit{Divergence} -- ситуация, когда вычислительные ядра одного варпа (warp) GPU выполняют сильно отличающийся код / обращаются к очень разным данным / etc
\pause
\item Самая важная оптимизация GPU path tracing'а -- уменьшить divergence
\end{itemize}
\end{frame}

\begin{frame}[fragile]
\frametitle{Divergence}
\usemintedstyle{solarized-light}
\begin{itemize}
\item Изначально каждый варп обрабатывает блок нескольких соседних пикселей
\pause
\item Что приводит к divergence?
\pause
\begin{itemize}
\item Какие-то лучи ушли в фон, а какие-то ещё нет
\pause
\item Обход BVH для разных лучей затронул разные части дерева
\pause
\item Разные лучи пересекли объекты с разными материалами и обращаются к разным текстурам
\pause
\item И т.п.
\end{itemize}
\end{itemize}
\end{frame}

\begin{frame}[fragile]
\frametitle{Wavefront path tracing}
\usemintedstyle{solarized-light}
\begin{itemize}
\item Один способ бороться с divergence -- разбить mega kernel на много небольших кусочков
\pause
\begin{itemize}
\item Ray generation -- генерирует луч для каждого пикселя и записывает их в большой буфер
\pause
\item Intersection -- вычисляет пересечения лучей со сценой
\pause
\item Material evaluation -- обрабатывает пересечения, генерирует новые лучи
\pause
\item И т.п.
\end{itemize}
\pause
\item Такой подход называется \textit{wavefront path tracing}
\end{itemize}
\end{frame}

\begin{frame}[fragile]
\frametitle{Wavefront path tracing}
\usemintedstyle{solarized-light}
\begin{itemize}
\item Много возможностей оптимизации:
\pause
\begin{itemize}
\item Как-нибудь отсортировать лучи перед шагом intersection, чтобы близкие в буфере лучи затрагивали близкие части BVH
\pause
\item Отсортировать пересечения по ID материала, чтобы близкие вызовы material evaluation обращались к одной и той же памяти
\pause
\item И т.п.
\end{itemize}
\end{itemize}
\end{frame}

\end{document}